% 02-muc-tieu.tex — Chương 2: Mục tiêu và phạm vi của đề tài

\chapter{Mục tiêu và phạm vi của đề tài}\label{ch:muctieu}

\section{Mục tiêu tổng quát của đề tài}

Triển khai các mô hình dự báo, dự đoán về hành vi của khách hàng, phục vụ công tác cải thiện chất lượng dịch vụ
và giữ chân người dùng. Báo cáo này tập trung vào bài toán dự báo khách hàng rời bỏ dịch vụ.

\section{Mục tiêu của bài toán: khách hàng rời bỏ}

Xây dựng và đưa vào vận hành một hệ thống tự động dự báo sớm khách hàng có nguy cơ rời bỏ dịch vụ, cung cấp cho đội chăm sóc khách hàng một danh sách ưu tiên kèm lý do rủi ro, qua đó hỗ trợ giữ chân khách hàng có giá trị và nâng cao hiệu quả điều hành.

\subsection{Mục tiêu cụ thể}

{\tablefont
\begin{longtable}{@{}C{0.8cm}L{4.4cm}>{\raggedright\arraybackslash}p{10.01cm}@{}}
\caption{Mục tiêu cụ thể của bài toán rời bỏ và kết quả đầu ra tương ứng}\label{tab:muctieu}\\
\toprule
\hdrrow\hh{TT} & \hh{Mục tiêu} & \hh{Kết quả đầu ra cần đạt}\\
\midrule
\endfirsthead
\toprule
\hdrrow\hh{TT} & \hh{Mục tiêu} & \hh{Kết quả đầu ra cần đạt}\\
\midrule
\endhead
\bottomrule
\endfoot
1 & Dự báo sớm & Mỗi kỳ, xác định được những khách hàng đang hoạt động có nguy cơ giảm hoặc ngừng sử dụng trong 1--2 tháng tới.\\
\rowa 2 & Tập trung nguồn lực & Xuất danh sách ưu tiên gồm nhóm 5\% khách hàng có mức rủi ro cao nhất, đủ ngắn để đội chăm sóc xử lý được.\\
3 & Giải thích được & Mỗi khách hàng trong danh sách có tối đa ba lý do chính, diễn đạt theo ngôn ngữ nghiệp vụ.\\
\rowa 4 & Tự động hóa & Dữ liệu được nạp, kiểm tra, xử lý và chấm điểm theo lịch cố định; không phụ thuộc thao tác thủ công.\\
5 & Duy trì chất lượng & Mô hình được huấn luyện lại định kỳ, chỉ thay thế mô hình cũ khi tốt hơn, và có giám sát biến động dữ liệu.\\
\rowa 6 & Đo được hiệu quả & Có khung chỉ tiêu đo lường và thiết kế thử nghiệm đối chứng để chứng minh giá trị kinh doanh.\\
\end{longtable}}
\vspace{4pt}

\section{Phạm vi của bài toán rời bỏ}

\subsection{Đối tượng và thời gian}

\begin{itemize}
  \item Đối tượng dự báo: từng khách hàng, được nhận diện bằng mã khách hàng đã mã hóa.
  \item Chân trời dự báo: hai tháng tiếp theo. Hệ thống nhìn vào hành vi các tháng gần nhất để dự báo tháng kế tiếp và tháng kế tiếp nữa.
  \item Độ dài lịch sử quan sát: tối thiểu ba tháng gần nhất; hệ thống tự thử nhiều độ dài khác nhau và chọn độ dài cho kết quả tốt nhất.
  \item Nhóm ưu tiên: khách hàng có sản lượng và doanh thu cao. Tài liệu nghiệp vụ nhắc đến ngưỡng khoảng 3 triệu đồng doanh thu mỗi tháng và nhóm gửi trên 100 đơn mỗi tháng; hệ thống đã có sẵn bộ lọc theo sản lượng và doanh thu để nâng ngưỡng khi nghiệp vụ yêu cầu.
\end{itemize}

\subsection{Nhóm khách hàng chưa đưa vào phiên bản này}

Để bảo đảm độ tin cậy, phiên bản hiện tại chủ động loại một số nhóm khách hàng mà việc dự báo dễ gây nhiễu:

\begin{itemize}
  \item Khách hàng mới: Khách hàng chỉ dùng dịch vụ một lần hoặc chưa đủ lịch sử hoạt động.
  \item Khách phụ hồi: Khách hàng quay lại sau thời gian ngừng.
  \item Khách mùa vụ: Khách hàng có tính mùa vụ rõ rệt, vì giảm đơn ngoài mùa là bình thường, không phải rời bỏ.
\end{itemize}

Ba nhóm này được đưa vào lộ trình phát triển với cách xử lý riêng, thay vì trộn lẫn làm sai lệch kết quả chính.

\section{Nguyên tắc thiết kế}

{\tablefont
\begin{longtable}{@{}>{\bfseries}L{4.0cm}>{\raggedright\arraybackslash}p{11.63cm}@{}}
\caption{Năm nguyên tắc thiết kế chi phối giải pháp}\label{tab:nguyentac}\\
\toprule
\hdrrow\hh{Nguyên tắc} & \hh{Ý nghĩa trong thực tế}\\
\midrule
\endfirsthead
\toprule
\hdrrow\hh{Nguyên tắc} & \hh{Ý nghĩa trong thực tế}\\
\midrule
\endhead
\bottomrule
\endfoot
Mỗi tầng một nhiệm vụ & Nạp dữ liệu, tạo chỉ số, huấn luyện, chấm điểm và giám sát được tách riêng; lỗi ở tầng nào xử lý ở tầng đó, không lan sang tầng khác.\\
\rowa Nhìn về phía trước, không ``nhìn trộm'' & Chỉ dùng thông tin có trước thời điểm dự báo để học; kiểm tra mô hình bằng cách giả lập đúng như đang vận hành tại từng thời điểm trong quá khứ.\\
Nhãn theo quy tắc nghiệp vụ & Nhãn rời bỏ được xác định bằng quy tắc rõ ràng, dựa trên đơn hàng thật, có tham khảo ý kiến đội chăm sóc khách hàng.\\
\rowa Không đánh đổi chất lượng lấy tốc độ & Mô hình mới chỉ được dùng khi vượt mô hình hiện tại và qua các kiểm tra an toàn.\\
Giải thích được & Kết quả luôn đi kèm lý do để con người có thể kiểm tra, phản biện và hành động.\\
\end{longtable}}
\vspace{4pt}
